\documentclass[11pt,a4paper]{article}
\usepackage[margin=23mm]{geometry}
\usepackage{amsmath,amssymb,array,booktabs}
\input{glyphtounicode}
\pdfgentounicode=1
\usepackage[hidelinks]{hyperref}
\hypersetup{hypertexnames=false}
\setlength{\parindent}{0pt}
\setlength{\parskip}{5pt}
\newcommand{\pos}[1]{[#1]_+}
\renewcommand{\labelitemi}{\(\bullet\)}
\begin{document}
\begin{center}
{\Large Selected S-E-P: clearance and a concrete recovery-source choice}\\[4pt]
15 October 2026 --- bounded analytical review, not a traffic experiment
\end{center}

\textbf{Status.} The user selects S-E-P's priority direction, closing S-E-A2.
The previously proposed cooperative and state/work formulas remain trial
laws. Two explicit recovery-source laws are formed below for comparison;
neither is selected. Stop at \textbf{SEP-R1}, Type 2 G3. Global hybrid
well-posedness, complete MP6 and Step 3 are not established.

\section*{1. Priority selection and the E-clearance boundary}

Retain $P_j=F_j(t-b/v)$, $q_j=P_j-Y_j$, $N_j=F_j-Y_j$, eligible
arrivals $\lambda_j=\dot P_j$, and initially empty terminal histories.
The independent B lane has $Y_B=P_B$. For the common triangular family,
\begin{equation}
 c(u)=\frac{wk_ju}{w+u},\quad C(v)=c(v),\quad
 U(f)=\frac{wf}{wk_j-f},\quad 0<R\leq C(v_-),\quad v\in[v_-,v_+]. \tag{1}
\end{equation}
Write $B(t)$ for shared available work, distinct from the B lane index.
S-E-P gives unqueued A first claim to $\lambda_A$ when it fits $B$.
Whenever A is queued, use the cooperative branch, including an empty E
offer. Thus cooperation does not disappear solely because E clears.

An exact local clearance diagnostic illustrates the distinction. Take
\begin{equation}
 B=4/5,\quad \lambda_A=3/5,\quad \lambda_E=1/10,\quad
 q_A>0,\quad q_E\downarrow0,\quad c_A=c_E=1,\quad C\geq B. \tag{2}
\end{equation}
The unit-cost limit is an accounting diagnostic, not a primary parameter
choice or an H1 reachability claim. Immediately before E clears, both
queues can receive $(g_A,g_E)=(2/5,2/5)$, so E drains. At E clearance,
the selected cooperative template serves
\begin{equation}
 (g_A,g_E)=(7/10,1/10),\qquad
 (\dot q_A,\dot q_E)=(-1/10,0). \tag{3}
\end{equation}
A literal ``strict A priority unless both queues are positive'' toggle
instead assigns $g_E=0$ at $q_E=0$, and $g_E=2/5$ at $q_E>0$.
It would require
\begin{equation}
 \dot q_E=1/10\ (q_E=0),\qquad
 \dot q_E=-3/10\ (q_E>0). \tag{4}
\end{equation}
There is no nonnegative absolutely continuous continuation from zero:
on any positive component $(a,b)$, integration gives
$q_E(t)=-3(t-a)/10<0$; staying identically zero contradicts (4).
An averaged/sliding rule would be another law, not the literal toggle.
Equation (3) is already the proposed S-E-P behavior. No new latch,
smoothing, positive queue floor or time-step convention is adopted.

\newpage
\section*{2. A speed-marked finite recovery profile}

Keep the prior local work family for conditional qualification,
\begin{equation}
 c_A=1,\qquad c_E=1+\delta+\beta\frac{g_A}{C(v)}
                      \frac{\pos{u_A-u_E}}v,\quad
 \delta,\beta\geq0,\quad g_A+c_Eg_E\leq B. \tag{5}
\end{equation}
The prior note proves fixed-regime unique queued roots under the sufficient
condition $1+\delta-\beta(1+v_+/w)>0$. Free trace marks use $u_j=v$;
queued trace marks use $u_j=U(g_j)$. Those representative speed and
current-work identifications remain hypotheses, not measured kinematics.

Let $a_r>0$ be a common physical recovery acceleration and
$m_{\rm ref}=1$ vehicle-equivalent fix normalization, not a packet size.
For an insertion-speed mark $u\geq0$, define
\begin{align}
 u_R&=U(R),&
 V(\theta;u)&=\sqrt{(w+u)^2+2wa_r\theta}-w,\nonumber\\
 D(\theta;u)&=\pos{R-c(V(\theta;u))},&\theta&\geq0. \tag{6}
\end{align}
Here $\theta$ is boundary-observed age, including backward propagation,
as in the earlier isolated recovery derivation. The positive part is the
proposed physical loss law, not a numerical clipping repair. The profile
is zero if $u\geq u_R$, and otherwise ends after
\begin{align}
 \Theta(u)&=\frac{u_R-u}{a_r}+\frac{u_R^2-u^2}{2a_rw},\qquad
 \Theta_{\max}=\Theta(0),\nonumber\\
 W(u)&=\int_0^\infty D(\theta;u)d\theta
   =\frac{k_c(u_R)\pos{u_R-u}^2}{2a_r},\qquad
 k_c(u)=\frac{wk_j}{w+u}. \tag{7}
\end{align}
Thus $0\leq D\leq R$, it decreases with age and speed mark, and
$W(u)\leq W_0:=W(0)<\infty$. On the active part, differentiating (6)
gives $|\partial_uD|\leq k_j$ and
$|\partial_\theta D|\leq a_rk_j$. These bounds extend across the
continuous zero-loss cutoff in the Lipschitz sense.

The isolated profile is a deduction using the point moving-bottleneck
case in Leclercq et al. (2016). Summing profiles from continuous mass is
an additional aggregate overlap assumption. The source's native wave/void
interactions are not reproduced exactly. In particular, it does not show
that every slow A departure should create a fresh independent profile.

\newpage
\section*{3. Two concrete laws for the same shared account}

Both candidates retain the same S-E-P allocation, trial work/trace laws,
empty initial history, common parameters and one budget output:
\begin{equation}
 B(t)=\pos{R-p(t)}. \tag{8}
\end{equation}
For each admitted stream, define a proposed speed-marked contribution
\begin{equation}
 p_j(t)=\frac1{m_{\rm ref}}\int_0^tD(t-s;u_j(s))\,dY_j(s)
       =\frac1{m_{\rm ref}}\int_0^t
           D(t-s;u_j(s))g_j(s)ds,\quad j=A,E. \tag{9}
\end{equation}

\textbf{R-E: E-origin disturbance, A follows the existing restriction.}
Set $p=p_E$. Actual admitted E quantity and its predicted speed determine
new impairment. A affects that speed and admission through the coupled
node and consumes the resulting shared service. Its slow discharge is
not separately a fresh source. This is not the old fixed-speed M1 kernel:
its mark varies with current predicted traffic state and actual admission.

\textbf{R-AE: both slow streams renew impairment.}
Set $p=p_E+p_A$. A passage with mark below $u_R$ also deposits a new
profile, even when its low speed resulted from a prior E restriction.
That is an explicit renewal/independence assumption. The kernel is neutral
for sufficiently fast passage, not triggered merely by a queue flag.

These are specified constitutive candidates once parameters and the trial
node are supplied, not names for unspecified future research. Their
physical speed marks, overlap and complete event proofs are unfinished.
R-AE's stronger assumption may count genuine successive acceleration
burdens, or may overcount followers already constrained by one disturbance.
R-E deliberately omits such independent follower renewal. Both limitations
require evidence; neither follows from conservation alone.

Internal source histories combine into one $p$ and one $B$. There is no
separate service ceiling per past passage, integer admission, hidden
initial speed oracle or longitudinal recovery zone. Only the last
$\Theta_{\max}$ of predicted admission/mark history can contribute.
This is recovery memory in the point account, not S-V approach-speed ODEs.

For a \emph{fixed} admissible mode schedule, the prior monotone inverse
makes $g_j$ and the trace marks locally Lipschitz in $B$. Equations
(6),(9) then give a causal Volterra map with difference bounded by
$K\Delta t\|\Delta B\|_\infty$ on a short interval with identical past.
Together with the 1-Lipschitz map in (8), this is a contraction for small
enough proof interval. No numerical time scale is introduced. It gives
local fixed-mode causal construction, not a theorem through arbitrary
simultaneous or grazing queue events or a full proof of the hybrid model.

\newpage
\section*{4. Exact recovery-tail contrast and finite-work bounds}

Suppose no E is admitted after $T_E$. Under R-E,
\begin{equation}
 p_E(t)=0\quad(t\geq T_E+\Theta_{\max}),\qquad
 \dot p_E\leq0\quad(t>T_E) \tag{10}
\end{equation}
almost everywhere. A can still be waiting during that recovery. Queue
clearance never resets $p$ artificially.

Consider an A-only queued tail, no new eligible arrivals, and
$0<B<R$. S-E-P gives $g_A=B$, and its trial queued speed mark is $U(B)$.
Because $c(U(B))=B$,
\begin{equation}
 D(0;U(B))=R-B. \tag{11}
\end{equation}
R-AE therefore adds new A loss at rate
\begin{equation}
 \frac{g_AD(0;u_A)}{m_{\rm ref}}
       =\frac{B(R-B)}{m_{\rm ref}}>0;\qquad
 R=1,\ B=3/4,\ m_{\rm ref}=1\ \Longrightarrow\ 3/16. \tag{12}
\end{equation}
R-E has no such new source. For identical inherited history at a comparison
instant, (12) is the difference in initial loss derivatives. Old profiles
also recover, so R-AE's total loss need not increase; it may recover more
slowly. No global ordering of two endogenously different trajectories is
claimed. The rational tail case is an exact local diagnostic, not a new
primary initialization or an H1 reachability certificate.

For either candidate, nonnegative kernels and mass conservation imply
by Tonelli
\begin{equation}
 \int_0^\infty p(t)dt\leq
 \begin{cases}
 M_EW_0/m_{\rm ref},&\text{R-E},\\
 (M_A+M_E)W_0/m_{\rm ref},&\text{R-AE}.
 \end{cases} \tag{13}
\end{equation}
$M_j$ denotes total mass eventually admitted from its terminal stream;
early-merged traffic belongs to A and is not forced to merge again.
The bound also holds before full clearance using the upper bound $M$.
Since $R-B=\min(R,p)\leq p$, (13) establishes a finite integrated service
loss for these particular laws, conditional on an admissible solution.

With $c_{\max}=1+\delta+\beta$, all mass eligible by $t_*=L/v$, and
$M_{\mathcal S}=M_E$ or $M_A+M_E$, the prior work-conserving bound now yields
\begin{equation}
 T_{\rm clear}\leq t_*+
 \frac{c_{\max}(q_A(t_*)+q_E(t_*))+M_{\mathcal S}W_0/m_{\rm ref}}R.
 \tag{14}
\end{equation}
Thus finite loss is no longer an unspecified assumption for either trial
source law. This does not prove existence, hybrid uniqueness or empirical
fidelity. No positive service floor or inherited M1 uniqueness proof is used.

\newpage
\section*{5. A limiting case that constrains the traffic claim}

Take an initially neutral account, empty A queue, constant A eligible rate
$\lambda_A=a<R$, and the selected free-A priority branch. While A remains
free, both recovery laws have no A source because $u_A=v\geq u_R$.
Using $c_0=1+\delta$, the instantaneous work law gives
\begin{equation}
 0\leq g_E\leq\frac{B-a}{c_0}. \tag{15}
\end{equation}
For the nonincreasing recovery profiles, differentiating (9) gives
\begin{equation}
 \dot p_E=\frac{g_E(t)D(0;u_E(t))}{m_{\rm ref}}+
 \frac1{m_{\rm ref}}\int_0^t
 \partial_\theta D(t-s;u_E(s))g_E(s)ds
 \leq\frac{R}{m_{\rm ref}}g_E. \tag{16}
\end{equation}
On this branch $B>0$ and $B=R-p_E$. Set $x=B-a$. Then
\begin{equation}
 \dot x\geq-\frac{R}{m_{\rm ref}c_0}x,\qquad
 x(t)\geq(R-a)\exp\!\left[-\frac{Rt}{m_{\rm ref}c_0}\right]>0.
 \tag{17}
\end{equation}
Hence A cannot begin queueing in finite time in this constant-A branch.
E may wait and reduce shared service, but its own admission is throttled
before the smooth loss crosses A's protected offer. Adding the R-AE source
does not create a first A queue there: its extra source is still neutral
until A actually becomes slow under the trial rule.

This is a specific limiting result, not an external-seed requirement or
a theorem that primary A delay is impossible. H1 produces time-varying
$\lambda_A=v(r-e(t-b/v))$ over its arrival pulse, generally increasing;
for varying $a(t)$ the inequality contains an additional $-\dot a$ term.
The constant-rate barrier therefore does not apply directly. A consequential
joined A-impact/MP6 witness remains to be established rather than assumed.

\textbf{Unchanged interfaces.} For the retained all-action domain
$M<k_jb$ and $\dot F_j\leq C(v)$, terminal vacancy stays positive and H1
is protected for any eventual admissible departures. Physical mass is
$N=F-Y$, and complete local TST remains
\begin{equation}
 J(v)=\frac{M(L-\ell/2)}v+\int_0^\infty(q_A+q_E+q_B)dt. \tag{18}
\end{equation}
Memory after completion is not vehicle-time. Scalar reconstruction and
gradient usefulness do not follow from these bounds; neither an interior
optimum nor full-lane-state necessity is required.

\textbf{Gate SEP-R1.} Do ordinary slow A departures merely follow the
represented disturbance (R-E), or renew it through new loss profiles
(R-AE)? Both are concrete source-law candidates; both need physical and
global-event qualification. Recommend R-E first as the minimal primary
candidate, not as proof independent A recovery is negligible. Neither source
law is selected. S-E-P stays selected and the whole review/Step 3 remains open.
\end{document}
